\section{Voraussetzungen}
\label{joukowski:sec:voraussetzungen}

Es handelt sich beim Satz von Kutta-Joukowski um ein approximatives mathematisches Modell. 
Damit wir dieses Modell für eine reale physikalische Situation anwenden können,
müssen einige Voraussetzungen gelten oder Annahmen zur Situation getroffen werden.
Nur wenn diese Voraussetzungen für die gegebene Situation gelten oder die Abweichungen davon vernachlässigbar klein sind,
kann das Modell angewendet werden.
Sonst kann es durchaus vorkommen,
dass das Resultat vom Modell eine grosse Diskrepanz zur Realität aufweist.

\subsection{Physikalische Voraussetzungen}

\subsubsection{Reibungsfreiheit}
	Reibungsfreiheit bedeutet,
	dass das Fluid,
	in unserem Fall die Luft, 
	keine Viskosität,
	d. h. innere Reibung,
	aufweist.
	Somit gibt es keine Tangential-, bzw. Scherkräfte,
	sondern nur Normal-, bzw. Druckkräfte auf die Oberfläche des Flügels.
	Zudem gilt die Bernoulli-Gleichung,
	welches unserem Modell zugrunde liegt,
	nur bei einer reibungsfreien Strömung.

\subsubsection{Inkompressibilität}
	Inkompressibilität bedeutet, dass die Dichte $\varrho$ des Fluides überall im Raum und zu jeder Zeit konstant ist.
	Somit gilt sie in allen Formeln nur als Skalierungsfaktor.
	Diese Annahme können wir bei unserer Luftströmung anwenden,
	sofern die Anströmungsgeschwindigkeit nicht zu gross wird.
	Unter $100\frac{\text{m}}{\text{s}}$ gilt diese Approximation sehr gut.
	Mathematisch formuliert können wir die Inkompressibilität als
	\begin{equation}
		\label{joukowski:voraussetzungen:inkompressibel}
		\nabla \cdot \vec{v}
		=
		0
	\end{equation}
	schreiben, wobei $\vec{v}$ die Geschwindigkeit des Fluides ist. 
	Diese Gleichung besagt, dass es nirgends im Raum eine Ansammlung des Fluides gibt. Die Dichte  bleibt somit überall im Raum konstant.

\subsubsection{Laminarität}
	Laminarität bedeutet, dass das Fluid in Schichten strömt, die sich nicht untereinander vermischen.
	Diese Voraussetzung führt direkt dazu, dass es in der Strömung keine Verwirbelungen gibt.
	Ohne Verwirbelungen gibt es keine Luftreibung, was bedeutet,
	dass in unserem Modell eine reine Auftriebskraft resultiert.
	
\subsubsection{Stationarität}
	Stationäre Strömung bedeutet kurz gesagt 
	``Alles sieht \emph{für alle Zeit gleich} aus'',
	aber es muss nicht überall im Raum gleich sein.
	Von Punkt zu Punkt im Strömungsfeld gibt es sehr wohl Unterschiede.
	
	In der Strömungslehre bedeutet die Stationarität,
	dass sich die physikalischen Strömungsgrössen wie Geschwindigkeit, Druck und Dichte an jedem Ort im Strömungsfeld mit der Zeit nicht ändern.
	
	Mathematisch formuliert bedeutet stationär
	\[
		\frac{\partial f}{\partial t}
		=
		0,
	\]
	wobei $f$ eine physikalische Strömungsgrösse darstellt.
	
	Die Rolle der Stationarität ist im Gegensatz zu den anderen physikalischen Voraussetzungen weniger offensichtlich.
	Bei einer stationären Strömung ist der Ausgleichsprozess abgeschlossen.
	Während einem Ausgleichsprozess kann sich die Strömung wesentlich anders verhalten, wie im ausgeglichenem Zustand.
	Dabei steigt die Komplexität so weit, dass sich die Strömung im Allgemeinen nur numerisch berechnen lässt.
	Dies beschreibt zum Beispiel die Situation beim Anfahren des Flugzeuges.
	Somit kann der Satz von Kutta-Joukowski nicht für diese Situation angewendet werden. 
	
\subsection{Situationsbeschreibung}
\label{joukowski:voraussetzungen:situationsbeschreibung}
\begin{figure}

	\centering
	\includegraphics[width=0.45\textwidth]{papers/joukowski/images/02_voraussetzungen/}
	\caption{
		
		\label{joukowski:voraussetzungen:fig_ball}
	}
\end{figure}

	Obwohl wir uns in einer dreidimensionalen Welt befinden,
	können wir uns in userem mathematischen Modell auf zwei Dimensionen begrenzen.
	Diese Vereinfachung ist nur möglich, da wir voraussetzen, dass die Strömung entlang des ganzen Flügels identisch ist.
	Somit müssen wir bei den Berechnungen nicht den gesamten dreidimensionalen Flügel betrachten,
	sondern können uns auf einen zweidimensionalen Querschnitt davon beschränken.
	Jetzt haben wir ein sogenanntes Flügelprofil, welches von einem inkompressiblen, reibungsfreien Fluid lamiar umströmt wird.
	In grosser Entfernung besitzt diese Strömung eine konstante Geschwindigkeit,
	welche wir mit $v_\infty$ bezeichnen. 
	Das $\infty$ bezieht sich darauf, dass in der Theorie dieser Punkt unendlich weit vom Flügelprofil entfernt sein müsste.
	Eine Folge der Entfernung einer Dimension ist,
	dass Grössen wie die Kraft $F$ zur Kraft pro Flügellänge $F'$ gewandelt werden.
	Wobei sich $F'$ mit
	\[
		F'
		=
		\frac{\partial F}{\partial l}
	\]
	berechnen lässt.
	
	Durch die Beschränkung auf zwei Dimensionen können wir zudem unsere Situation mittels der komplexen Zahlenebene vollständig beschreiben.
	In den folgenden Abschnitten wird also die örtliche Abhängigkeit mit $z = x + iy$ beschrieben. Manchmal wird die vertikale Koordinate $y$ auch mit $h(z)$ bezeichnet,
	weil sie zum Beispiel die Höhe eines fliegenden Flugzeugs darstellt. $h_\infty$ beschreibt die Höhe im Referenzpunkt.
	Des weiteren können wir auch zweidimensionale Vektoren als komplexen Zahlen darstellen.
	Dies kann bei der Geschwindigkeit $\vec{v}$ als
	\[
		\vec{v}
		=
		v(z)
		=
		v_x + i v_y
	\]
	geschrieben werden.
	Dabei ist es Konvention,
	\[
		v(z)
		=
		u + iw
	\]
	zu schreiben.
	Dabei entspricht $u$ der sogenannten Längskomponente und $w$ der sogenannten Querkomponente der Geschwindigkeit $v(z)$.
	
	Dadurch, dass wir unsere Grössen mit den komplexen Zahlen formuliert haben, können wir alle Werkzeuge, welche uns die komplexe Analysis zur Verfügung stellt, für den Rest unseres Kapitels nutzen.